\chapter{Limits and Homing}\label{ch:limits}

Limit switches and homing are optional. Without them, the board does not know
where the ends of the axes are: the plot area and the origin are all that keep
the pen on the bed (section~\ref{sec:plot-area}). With them:

\begin{itemize}
  \item A \term{limit switch} stops the motion of an axis toward the end of
        its travel. A job that reaches a limit stops, and the board drops it.
  \item A \term{home switch} gives the axis a known position. After homing,
        plotter 0,0 is at the same place on the bed each time.
\end{itemize}

The switches belong to the axes as HP-GL sees them (X and Y), not to the
motors. Thus they operate the same way on a Cartesian machine, on a CoreXY
machine, and with \key{dir_x}, \key{dir_y}, or \key{swap_xy}.

\begin{warning}
A limit switch is not an emergency stop. The board stops the axis with a
ramp, so the axis moves past the switch before it stops. Connect an emergency
stop to \conn{DI-6} (section~\ref{sec:fault-behavior}) and keep space for the stop
distance after each limit switch.
\end{warning}

\section{How the Board Uses the Switches}\label{sec:switch-rules}

\paragraph{Inputs.} Each switch uses one input of the ClearCore:
\conn{DI-7}, \conn{DI-8}, \conn{A-9} to \conn{A-12}, or \conn{IO-0} to
\conn{IO-5}. The board does not let two switches use the same input, with one
exception: the home switch of an axis can use the input of a limit switch of
the same axis. Then the axis homes on that limit switch. \conn{DI-6} is for the
emergency stop, and the pen output (\key{pen_io}) cannot be a switch.

\paragraph{Contacts.} A \term{normally closed} switch is on at rest and off
when it trips. This is the default and the recommended contact: a broken wire
reads as a tripped switch, so the axis stops. A \term{normally open} switch is
off at rest. Set its bit in \key{switch_open} (table~\ref{tab:set-switches}).

\paragraph{Filter.} The board reads the inputs at 5~kHz, in the same interrupt
that makes the steps. A switch changes state when its input has the new level
for three samples in a row (0.6~ms). The board keeps the position of the first
of the three samples. This is the position where the switch tripped.

\paragraph{Limits.} A tripped limit switch stops only the motion
\emph{toward} it. Motion away from it is always possible, so you can move the
axis off the switch. The stop is a ramp along the path at three times
\key{accel_mm_s2}. It starts in the same sample in which the switch trips. The
distance after the switch is approximately:
\begin{equation}
  d = \frac{v^{2}}{2 \times 3 \times a} + v \times 0.6\ \text{ms}
\end{equation}
where $v$ is the speed and $a$ is \key{accel_mm_s2}. For example, at
200~mm/s and 2000~mm/s\textsuperscript{2}, $d \approx 3.5$~mm. At a jog speed
of 10~mm/s, $d$ is less than 0.02~mm.

\begin{table}[htbp]
\centering
\caption{What a limit switch does}\label{tab:limit-actions}
\begin{tabularx}{\linewidth}{@{}P{0.28\linewidth}L@{}}
\toprule
\textbf{Motion} & \textbf{Result} \\
\midrule
A job, \cmd{goto}, or \cmd{pen} & The axis stops. The board drops the job and
raises the pen. The fault \code{limit X+ (DI-7)} (for example) latches. The
motors stay enabled. The next job waits until you type \cmd{clear}. \\
A jog & The jog stops. This is not a fault. A jog toward the tripped limit is
refused; a jog away from it is accepted, and clears the fault. \\
Homing & Homing stops with the error \code{limit X- reached before the home
switch} (for example), unless the home switch uses the input of that limit. \\
\bottomrule
\end{tabularx}
\end{table}

\section{Connect the Switches}

\begin{caution}
Disconnect the power of the ClearCore before you connect the switches.
\end{caution}

\begin{procedure}{Connect a limit or home switch}
\begin{enumerate}
  \item Use a switch with a contact for the input voltage of the ClearCore
        (24~V DC). Use normally closed contacts if possible.
  \item Connect the switch between the 24~V supply and the input, as the
        ClearCore manual shows for a digital input.
  \item Put the limit switch where the carriage trips it before the end of the
        travel. Keep the stop distance $d$ (section~\ref{sec:switch-rules})
        between the switch and the end.
  \item Write down the input of each switch. The wizard finds them again, so
        a mistake here is not a problem.
\end{enumerate}
\end{procedure}

\section{Set Up the Switches With the Wizard}\label{sec:switch-wizard}

The desktop app has a wizard that finds each switch, reads its contact, and
proves each one by motion. It sets one axis at a time. Do it again for the
other axis.

\begin{procedure}{Set up the switches of an axis}
\begin{enumerate}
  \item Connect the app to the plotter (section~\ref{sec:app-connect}).
  \item Select \menu{Plotter\,>\,Set up limits and homing\ldots}.
  \item \menu{Axis}: select X or Y and the switches that the axis has. If the
        axis homes on a limit switch, select it in \menu{Home switch is}.
        Keep the setup jog speed low (10~mm/s). Click \menu{Next}.
  \item \menu{Find each switch}: click \menu{Find} for a switch, then trip
        the switch by hand for approximately one second and release it. The
        wizard shows the input that changed. Do this for each switch.
  \item \menu{Contact}: make sure that the contact of each switch is correct.
        The wizard reads it from the level at rest.
  \item \menu{Direction}: click \menu{Enable motors}. Hold \menu{Hold: jog +}
        until a switch stops the axis, then hold \menu{Hold: jog -}. The
        plotter stops at the first switch of the axis. If the axis goes
        toward + and stops on the - limit, the wizard asks a question:
        \begin{itemize}
          \item If the head went to the - end, click \menu{The head went the
                other way: reverse the axis}. The wizard changes \key{dir_x}
                or \key{dir_y} and disables the motors. Enable them and jog
                again.
          \item If the head went to the + end, click \menu{I tripped them the
                wrong way round: swap them}.
        \end{itemize}
  \item \menu{Limits}: the wizard writes the limits to the plotter. Hold the
        jog until each limit stops the axis, then jog back off it. The wizard
        shows how far the axis moved past each switch.
  \item \menu{Home}: make sure that the home direction is correct. Set the
        position of the axis at the switch: 0 for a switch at the 0 end of the
        plot area, or the width (or height) for a switch at the far end. Click
        \menu{Write and home}. The plotter homes the axis the number of times
        that you set, and the wizard shows the spread of the positions.
  \item \menu{Review}: select \menu{Home before plotting} if the plotter must
        home before each job after power-up. Click \menu{Apply}.
\end{enumerate}
\end{procedure}

\begin{note}
\notelabel The Limits and Home steps write their settings to the plotter
immediately. \menu{Close} does not remove them. To remove a switch, set its
input to \code{none} in \menu{Plotter settings}.
\end{note}

\section{Set Up the Switches in the Console}

You can also set the switches with \cmd{set}. This example has a normally
closed X- limit on \conn{DI-7}, which is also the X home switch, and a
normally closed Y- limit and home switch on \conn{DI-8}:

\begin{lstlisting}[language=console]
set x_limit_minus DI-7
set x_home DI-7
set y_limit_minus DI-8
set y_home DI-8
set x_home_dir -1
set y_home_dir -1
set home_mode 1
home
\end{lstlisting}

Use \cmd{swinfo} to see the levels of the inputs, the state of each switch,
the limit stops, and the last homing (appendix~\ref{app:console}).

\section{Homing}\label{sec:homing}

\cmd{home} (or \menu{Home} in the app) homes each axis that has a home switch.
\cmd{home x} or \cmd{home y} homes one axis. For each axis, in the sequence
that \key{home_order} sets, the board:

\begin{enumerate}
  \item Raises the pen, if it is down.
  \item Moves off the home switch, if the switch is tripped.
  \item Moves toward the switch at \key{x_home_fast} until the switch trips.
        If it does not trip within \key{x_home_travel} (or 1.2 times the
        plot area plus 20~mm), homing stops with an error.
  \item Moves back until the switch releases, then \key{x_home_back} more.
  \item Moves toward the switch at \key{x_home_slow} until it trips again.
  \item Sets the position where the switch tripped to \key{x_home_at}. The
        position is a whole number of steps, so each homing gives the same
        result.
\end{enumerate}

The head stops a small distance past the switch. Homing lasts for each axis
approximately the search distance divided by the search speed, plus a few
seconds.

\begin{procedure}{Home before each job}
\begin{enumerate}
  \item Set up homing (section~\ref{sec:switch-wizard}).
  \item Type \cmd{set home_required 1}, or select \menu{Home before
        plotting} in the wizard.
  \item Send a job. If the axes are not homed, the board homes them first and
        then plots. While it homes, the state is \code{homing}.
\end{enumerate}
\end{procedure}

The board forgets the homed position when you disable the motors, when the
emergency stop opens, when a servo has a fault, and when you set the origin.
The servos can move when they are disabled, so the position is not sure after
that.

\paragraph{ClearPath hard-stop homing.} With \key{home_mode}~2, the ClearPath
servos home themselves against a hard stop, with no switch. This is for
Cartesian machines only, because each servo homes one axis.

\begin{procedure}{Use ClearPath hard-stop homing}
\begin{enumerate}
  \item In ClearPath MSP, set homing for each servo of X and Y. Set it to home
        on every enable.
  \item Type \cmd{set hlfb_check 1}. The board sees the end of homing on HLFB.
  \item Type \cmd{set x_home_at 0} and \cmd{set y_home_at 0}, or the positions
        where the hard stops are.
  \item Type \cmd{set home_mode 2}.
  \item Type \cmd{home}. The board disables the servos, enables them again,
        and waits up to 120~seconds for HLFB. Then it sets the position.
\end{enumerate}
\end{procedure}

\section{Homing Errors}

\begin{table}[htbp]
\centering
\caption{Homing errors}\label{tab:homing-errors}
\begin{tabularx}{\linewidth}{@{}P{0.42\linewidth}L@{}}
\toprule
\textbf{Message} & \textbf{Action} \\
\midrule
\code{X home switch not found within the search travel} & Make sure that the
switch operates (\cmd{swinfo}) and that \key{x_home_dir} goes toward it.
Increase \key{x_home_travel}. \\
\code{limit X- reached before the home switch} & The home switch is beyond the
limit, or \key{x_home_dir} is wrong. \\
\code{X home switch did not release} & The switch stays tripped. Examine the
switch and its wire. Increase \key{x_home_back}. \\
\code{X home switch lost on the slow approach} & The switch did not trip again.
Increase \key{x_home_back}, or examine the switch for bounce. \\
\code{no X home switch (set x_home)} & \cmd{home x} needs a home switch. \\
\code{ClearPath homing did not finish (HLFB not asserted)} & Examine the homing
settings in MSP and the HLFB cable. \\
\bottomrule
\end{tabularx}
\end{table}

\section{Test Without Switches}\label{sec:sim-switches}

\cmd{sim} simulates a switch on an input, for a test on the bench. A simulated
switch is fixed to the machine: it stays at the same place when homing or
\cmd{origin} changes the position.

\begin{procedure}{Test homing with a simulated switch}
\begin{enumerate}
  \item Type \cmd{set x_home A-9}, \cmd{set x_home_dir -1}, and
        \cmd{set home_mode 1}.
  \item Type \cmd{sim A-9 x- 30}. A switch is now 30~mm from the head, on the
        - side of X.
  \item Type \cmd{enable}, then \cmd{home x}.
  \item Type \cmd{swinfo}. Make sure that the \code{home} line shows
        \code{homed=1} and \code{last=ok}.
  \item Type \cmd{sim clear}. Set the settings back.
\end{enumerate}
\end{procedure}
